\section{Calibrating the equivalence at a regular ray}\label{sec:raycalibration}

At the regular tuple $a$ the periods of the Laurent polynomial are known. We must compare them with a functor whose source is the particular stopped category $\mathcal C_a$. We use a second path for this purpose, joining $a$ to $*$ by interpolation of the \emph{real} logarithms of the ray coefficients. No vanishing path is transported through a collision on this second path. The quantities we follow are the far fibre and the cycles associated with ample line bundles. They remain defined when critical points in the interior coalesce.

\subsection{The class of a geometric brane}
For a compact Liouville domain $Y$ with contact boundary $C$ and closed stop $Q\subset C$, set $H(Q)=H_n(Y,C\setminus Q;\Z)$. In the rest of this section we work without the group, fix an ordered lattice frame, and transport the corresponding square root of the squared volume and the frame of the background bundle. The grading orients a rank-one disk: on a positively oriented frame the phase of the volume is $e^{\pi i\vartheta}$ for the chosen lift $\vartheta$ of the squared phase. Shifting the grading by one reverses the orientation. Denote by $\ell_Q(L)$ the relative class of its oriented truncation, and use the same notation for any similarly decorated cylindrical Lagrangian with trivial rank-one coefficients.

First identify the rightward relative group at a ray tuple $c$. Take a regular value $w_\delta=Re^{i\delta}$ for small $\delta>0$, and choose $q<R\cos\delta$ above the real parts of all critical values and compact perturbation images. Denote the compact page at $w_\delta$ by $P_\delta$, the truncated marked page at $R$ by $F^-$, and the full mirror torus by $T$. The comparison $J_c$ is the following pair of inclusions, read from left to right:
\begin{equation}\label{eq:relativejunction}
 H(F^-)\ \xleftarrow{\ \sim\ }\ H_n(Y,P_\delta;\Z)\
       \xrightarrow{\ \sim\ } H_n(T,\{\operatorname{Re} W_c>q\};\Z).
\end{equation}
For the first arrow, the vertical boundary is the mapping torus of the fibre and the horizontal boundary is the product of its frontier with a disk. Cutting the vertical boundary along a full page away from its outer collar leaves the inner fibre times an open interval. In the remaining binding collar, the base and collar coordinates form a disk with a radial slit running in from its outer edge to the truncated page frontier. This region and its overlap with the inner region both retract fibrewise to the binding: the transverse factors are the slit disk and a product of two open intervals. The homotopy pushout therefore identifies the inclusion of $P_\delta$ in the page complement as a homotopy equivalence. One may move the truncation frontier through regular collar levels before this calculation, round the corner away from a slightly smaller $P_\delta$, and restore its collar afterwards. With empty binding one simply cuts the mapping torus. For the second arrow, $Y\hookrightarrow T$ is a homotopy equivalence by completion at the central ray. Along the positive coefficient path the rounded domains are smoothly isotopic in $T$ by extension of their compact family of embeddings. Over the half-plane in \eqref{eq:relativejunction} horizontal lifting along straight segments to $w_\delta$ is complete on each finite segment by Lemma~\ref{lem:rayest}. This retracts the half-plane preimage onto its full fibre, whose compact truncation has only a collar removed. The sequence of a pair proves both isomorphisms. In the moving horizontal truncations the boundary product used here is the one obtained from exterior parallel transport in Proposition~\ref{prop:compactpair}; isotoping regular horizontal levels in this exterior gives the same page inclusions. A thimble ending at another small positive angle may be capped to $P_\delta$ within the outer circle. Both the cap and the outgoing tail lie in the indicated right half-plane for a sufficiently small angular interval.

\begin{lemma}\label{lem:braneclass}
On $\mathcal C_a$ there is an integral isomorphism $c_a:K_0(\mathcal C_a)\to H(F^-)$ sending every geometric brane as above to its oriented class. Along a deformation or change of section as in Lemma~\ref{lem:geometrictransport}, if $T_*$ is the categorical equivalence and $T_H$ is ordinary relative transport, the map $T_Hc_a K_0(T_*)^{-1}$ on the endpoint likewise gives the fundamental class of every such brane. The same holds on replacing a ribbon stop by its core.
\end{lemma}

\begin{proof}
At the uncut regular ray use the exact Lefschetz presentation with the marked page. Cut the open book at this page to obtain a Liouville sector. Its free ideal boundary is a contact neighbourhood of a complementary page up to deformation: the vertical open book cut at one angle is trivial over an interval and joins the product binding chart. This is the Lefschetz example of \cite[Examples~1.4, 2.19--2.20]{GPSI}. The resulting cutoff Reeb field has no periodic orbit by \cite[Lemma~2.36]{GPSI}; the acceleration $H^*(Y_{\rm sec},\partial Y_{\rm sec})\to SH^*(Y_{\rm sec},\partial Y_{\rm sec})$ is an isomorphism by \cite[Proposition~4.42]{GPSI}.

Apply the categorical Chern character to a perfect brane, then the open--closed map and inverse acceleration. Take its degree-$n$ component and use Poincaré--Lefschetz duality. In the compact sector the face relative in cohomology is complementary to the free face where our truncations end. The result is a homomorphism to $H(F^-)\otimes\C$. For the grading conventions of \cite{GPSI}, fold degrees modulo two before the trace and project from parity $n$ after inverse acceleration. The ordinary background frame supplies the Spin structure from the graded oriented relative Pin brane. Folding preserves determinant lines, signs of rigid counts and their parity; it extends to twisted complexes and retracts. The counts over orientation lines extend to $\C$. Morita invariance gives the trace on the perfect envelope. On a geometric brane its identity class goes to the acceleration of its Gysin class by \cite[Proposition~5.13]{GPSI}; after duality this is precisely $\ell(L)$. The sector-to-stopped inclusion carries the brane to the same brane with its collar. In the reverse direction an admissible brane can be moved into the sector by an isotopy avoiding the stop \cite[Lemma~3.8]{GPS}, with the same relative chain class.

The dragged distinguished disks at $a$ form a full exceptional collection by Lemma~\ref{lem:regularpath}. Their integral classes form a basis by triangularity and fullness. Their thimbles form the integral relative handle basis over a regular page, hence also a basis in \eqref{eq:relativejunction}. The map just constructed therefore lands exactly in the integral lattice and is an isomorphism there. This argument leaves the grading and coefficients of the original Grothendieck group in place.

We verify the assertion for inverse arrows as well. In the fixed contact boundary let $A\subset B$ be convex thickenings with one inward contact flow $\psi$ satisfying $\psi^\epsilon B\subset A$. For a brane $L$ avoiding $A$, lift $\psi^{-u}$ near infinity by its homogeneous Hamiltonian cut off towards the interior. Its image $L_u$ still avoids $A$, and $L_\epsilon$ avoids $B$. Exact isotopy identifies the two representatives in the $A$-stopped category. Its trace has side in $C\setminus A$, so the inclusion $H(B)\to H(A)$ sends $\ell_B(L_\epsilon)$ to $\ell_A(L)$. This homology inclusion is an isomorphism: in the nested sequence $B\supset A\supset\psi^\epsilon B\supset\psi^\epsilon A$ the shrink composites have homotopy inverses given by contact expansion extended over a collar. These same maps implement ordinary transport for nearby moving thickenings. A slightly enlarged common thickening contains them all, and a uniform finite contraction lies inside all of them. A chain allowed for the enlargement stays allowed throughout the short interval and gives the transport of its class. Subdivide the interval as in Lemma~\ref{lem:geometrictransport}.

Finite inward contractions of a thickening decrease to its core. Their open complements exhaust the core complement, so any finite chain or homology between finite chains occurs at a finite stage. The inclusion into the core-complement relative group is thus an isomorphism. It factors through the ribbon-complement group. For surjectivity to this intermediate group expand the boundary of a chain avoiding the ribbon by the inverse flow until it avoids the thickening. The negative contraction sends the ribbon into itself, so this trace stays in the ribbon complement. The same representative argument for branes uses a compact link disjoint from the core and one finite shrink missing that link. Thus the relative map agrees on the actual branes in both directions at all junctions; no transport of a singular core by a diffeomorphism is needed.
\end{proof}

Follow the compact pairs with the uncut $W_c$ on the positive path from $a$ to $*$. The face estimates of Lemma~\ref{lem:regularpath} concern only the regular boundary annuli and hold even if the interior has a degenerate critical point. The total native form is Kähler there since $W_c$ is holomorphic. At $*$ carry the pair to the positive conic link by Lemma~\ref{lem:equivariantlink}. At this endpoint reflect momentum, writing $m=\theta/(2\pi)$ and $y=-2\pi p$, so the primitive becomes $-y\,dm$ and the conormal cones in $y$ are the negative fan cones. This is a strict change of coordinates. Denote the composite equivalence by $T_*$, and the transport of relative homology by $T_H$. To calibrate it, let $P$ run over integral ample moment polytopes and write $C_P$ for the signed cycle of \eqref{eq:cycledual}. We shall calculate $T_H J_a^{-1}[C_P]$ on the cotangent end and realize the answer by an actual disk.
